\uuid{2wbx}
\titre{Géométrie}
\auteur{bodin ; hanani ; mzari}
\organisation{exo7}
\licence{CC BY-NC-SA 4.0}
\datecreate{2020-01-20}
\niveau{L1}
\module{Algèbre}
\chapitre{Nombres complexes}
\sousChapitre{Racine carrée, équation du second degré}
\difficulte{3}
\theme{Équations, géométrie}
\exo7qcmid{lille-1.nombres-complexes.029}

\contenu{
\begin{question}
\qclassification{104.02,104.03,104.04}
\qtags{motcle=géométrie}
Soit $E$ l'ensemble des points $M$ d'affixe $z$ tels que : $z+\frac{1}{z} \in \Rr$. Quelles sont les assertions vraies ?
\begin{answers}
    \bad{$E = \Rr^*$}

    \bad{$E$ est le cercle unité.}

    \good{$ E = \Rr^* \cup \{z\in \Cc; \, |z|=1\}$}

    \good{$E$ contient le cercle unité.}

\end{answers}
\begin{explanations}
Soit $z \neq 0$. On a :  $z+\frac{1}{z} \in \Rr \Leftrightarrow z+\frac{1}{z}  = \overline{z}+\frac{1}{\overline{z}}$. On multiplie par $z\overline{z}$ et on simplifie cette égalité.

\end{explanations}
\end{question}
}
